% Generated draft: requires mathematical and source review.
\documentclass[11pt]{article}
\usepackage[margin=25mm]{geometry}
\usepackage{fontspec}
\setmainfont{Latin Modern Roman}
\usepackage{amsmath,amssymb,mathtools,unicode-math}
\setmathfont{Latin Modern Math}
\usepackage{enumitem,needspace,xurl,hyperref}
\setlength{\parindent}{0pt}
\setlength{\parskip}{5pt}
\setlength{\emergencystretch}{3em}
\newcommand{\N}{\mathbb{N}}
\newcommand{\Z}{\mathbb{Z}}
\newcommand{\Q}{\mathbb{Q}}
\newcommand{\R}{\mathbb{R}}
\newcommand{\C}{\mathbb{C}}
\newcommand{\F}{\mathbb{F}}
\newcommand{\A}{\mathbb{A}}
\newcommand{\PP}{\mathbb{P}}
\newcommand{\E}{\mathbb{E}}
\newcommand{\eps}{\varepsilon}
\newcommand{\dd}{\,\mathrm{d}}
\begin{document}
\section*{UnsolvedMath ID: 30005173}
\textbf{Source problem code:} catalogue entry
\subsection{Definitions and mathematical statement}
A compact Kähler Poisson manifold is a pair \((X,\pi)\), where \(X\) is a compact complex manifold admitting a Kähler form and \(\pi\in H^0(X,\bigwedge^2T_X)\) is a holomorphic bivector satisfying \([\pi,\pi]=0\), with the bracket here denoting the Schouten bracket. It defines a Poisson bracket on local holomorphic functions by \(\{f,g\}=\pi(\dd f,\dd g)\). Write \(\pi^{\sharp}:T_X^*\to T_X\) for the induced bundle map. A symplectic leaf \(L\) is a connected maximal immersed complex submanifold whose tangent spaces are the images of \(\pi^{\sharp}\); its induced holomorphic Poisson structure is nondegenerate. Compactness of a leaf refers to its intrinsic manifold topology.

\textbf{Problem 6, as supplied.} Describe the compact Kähler Poisson manifolds \((X,\pi)\) for which there exists a compact symplectic leaf \(L\subseteq X\) such that
\[
\lvert\pi_1(L,\ell)\rvert=\infty
\]
for a base point \(\ell\in L\). Since \(L\) is connected, infinitude is independent of the base point. The source does not specify an equivalence relation for this description or a precise form of the desired classification.

\textbf{Supplied reduction.} The report passage asserts that this problem reduces to the case where \(L\) is a finite étale quotient of a torus. In the holomorphic setting, this means that there is a finite surjective étale holomorphic map \(T\to L\), where \(T=\C^g/\Lambda\) is a compact complex torus and \(\Lambda\) is a lattice of rank \(2g\). The infinite-fundamental-group condition requires \(g\geq 1\). This reduction is recorded as a source assertion; its proof and the precise manner in which it recovers the general case are not supplied.

\textbf{Qualifications from the surrounding passage.} The supplied Corollary 4 states that if \((X,\pi)\) is compact Kähler and
\[
h^{2,0}(X):=\dim_{\C}H^0(X,\Omega_X^2)=0,
\]
then every compact symplectic leaf is a point. The hypothesis \(h^{2,0}(X)=0\) belongs to that corollary and is not imposed in Problem 6. The preceding definition calls a holomorphic map \(\phi:(X,\pi_X)\to(Y,\pi_Y)\) a Poisson morphism when \(\phi^*\{f,g\}_Y=\{\phi^*f,\phi^*g\}_X\) for every pair of local holomorphic functions \(f,g\) on \(Y\).

The supplied Remark 5 says that the conclusion of an earlier theorem may fail if the complex Poisson manifold \(X\) is noncompact, if \(X\) is non-Kähler, or if the fundamental group of \(L\) is infinite. It also reports smooth real Poisson examples in which both the ambient manifold and its leaf are simply connected but the ambient manifold does not decompose as a product. The earlier theorem and the intended product factors are absent from the record, so these contextual assertions do not determine an additional precise theorem or a product-decomposition requirement for this problem.
\subsection{Short English statement}
Describe compact Kähler manifolds with a holomorphic Poisson structure that has a compact symplectic leaf with infinite fundamental group. The supplied report asserts a reduction to leaves that are finite étale quotients of compact complex tori. The intended form of the description and the details of this reduction remain unspecified.
\subsection{Sources}
\begin{itemize}
\item Supplied local problem registry record, permanent UnsolvedMath ID \(30005173\), exact source code \texttt{catalogue entry}, title \emph{Kähler Poisson Manifolds with Infinite-Fundamental-Group Leaves}. The record identifies the extracted question as Problem 6 and gives the proposed year as 2022; no individual proposer is supplied.
\item Supplied source citation: \emph{Algebraic Geometry: Moduli Spaces, Birational Geometry and Derived Aspects} (2023), Oberwolfach Reports, pp. 1805--1864, DOI \texttt{10.4171/owr/2022/32}. Supplied source URL: \url{https://doi.org/10.4171/owr/2022/32}. The record includes surrounding text containing Corollary 4, Remark 5, Problem 6, and the asserted torus-quotient reduction.
\item Supplied literature source: \emph{Oberwolfach Report 32/2022}, \url{https://doi.org/10.4171/owr/2022/32}. The literature assessment dated \texttt{2026-08-22} labels the problem open and describes it as the report's remaining Problem 6. This is supplied dated curation, not an independently verified assessment.
\item The registry labels the statement \texttt{corrected\_verified} and reports that the official report was checked and its notation repaired. These are supplied verification claims; no independent source inspection or mathematical verification is asserted here.
\end{itemize}
\end{document}
